% Computer Science A --- Gradle build automation (practice bank).
% Moved from \texttt{cs-web-adv/test.tex}; \texttt{cs-web-adv/} module.

\runningheader{Computer Science A}{Build automation (Gradle)}{Page \thepage\ of \numpages}

\begingradingrange{csabuildauto}

\uplevel{\par\textbf{Gradle}\par}

\question[4]
Which repositories can Gradle use for resolving dependencies? \emph{Select all that apply.}
\begin{checkboxes}
  \CorrectChoice Maven Central repository
  \choice Docker Hub
  \CorrectChoice Ivy repositories
  \CorrectChoice Local repositories
\end{checkboxes}
\begin{solution}
  Gradle resolves artifacts from Maven-compatible repos (including Maven Central), Ivy repos, and the local cache (typically \lstinline[style=bashinline]|~/.gradle/caches| or a declared local directory). Docker Hub distributes container images, not Java library coordinates for Gradle dependency resolution.
\begin{lstlisting}[style=java]
repositories {
    mavenCentral()
    ivy { url "https://example.com/ivy" }
    mavenLocal()
}
\end{lstlisting}
\end{solution}

\question[4]
For Gradle projects, a unit test library like JUnit would go in which section?
\begin{choices}
  \choice \texttt{plugins}
  \choice \texttt{repositories}
  \CorrectChoice \texttt{dependencies}
  \choice \texttt{tasks}
\end{choices}
\begin{solution}
  Library coordinates (for example \lstinline[style=javainline]|testImplementation 'org.junit.jupiter:junit-jupiter:5.10.0'|) belong in the \textbf{dependencies} block. \texttt{plugins} apply Gradle plugins; \texttt{repositories} declare where to download artifacts; \texttt{tasks} define or configure build steps.
\begin{lstlisting}[style=java]
dependencies {
    testImplementation 'org.junit.jupiter:junit-jupiter:5.10.0'
}
\end{lstlisting}
\end{solution}

\question[4]
What are Gradle build scripts?
\begin{choices}
  \CorrectChoice Script files that define Gradle \textbf{projects} and \textbf{tasks}; a project may represent a library JAR, web application, or assembled deliverable
  \choice Script files that list dependencies stored as JARs in the project root package only
  \choice XML project descriptors with default directories such as \texttt{src/main/java} (Maven \texttt{pom.xml} style)
  \choice JVM instruction files compiled with \lstinline[style=javainline]|javac| into \texttt{.class} files
\end{choices}
\begin{solution}
  Gradle builds are driven by scripts (commonly \lstinline[style=bashinline]|build.gradle| or \lstinline[style=bashinline]|build.gradle.kts|). Each build has one or more \textbf{projects}; each project registers \textbf{tasks} (compile, test, jar, etc.). Dependencies are declared in the script but downloaded to Gradle's cache, not checked into the repo root. The XML description matches Maven, not Gradle; build scripts are not Java source compiled by \lstinline[style=javainline]|javac|.
\end{solution}

\question[4]
What is Gradle?
\begin{choices}
  \choice A dependency management tool only
  \choice The Java compiler
  \choice A container orchestrator
  \CorrectChoice A build automation tool
\end{choices}
\begin{solution}
  \textbf{Gradle} automates compiling, testing, packaging, and publishing software. It \emph{includes} dependency management but also orchestrates tasks and plugins across projects. The Java compiler is part of the JDK; container orchestration (for example Kubernetes) is a separate concern.
\begin{lstlisting}[style=bash]
./gradlew build    # compile, test, and assemble artifacts
./gradlew test     # run unit tests
\end{lstlisting}
\end{solution}

\endgradingrange{csabuildauto}
